\begin{titlepage}

\vspace*{0.75in}

\begin{center}

\textbf{\sffamily\bfseries\huge
Systems Biology Graphical Notation:\\[0.3em]
\PDl Level 1}

\vspace*{0.5in}

\Large
\textbf{Version 2.1}\\[0.1in]

\large
\today\\[0.25in]

\vspace{0.5in}

\textbf{\sffamily Editors}:\\[7pt]

\begin{tabular}{l>{\hspace*{15pt}}r}
Hasan Balci   & \emph{NIH, USA}\\
Adrien Rougny   & \emph{University of Luxembourg, Luxembourg}\\
Rupert Overall   & \emph{Humboldt University, Germany}\\
Irina Balaur & \emph{EISBM, France}\\
Michael Blinov & \emph{UConn School of Medicine, USA}\\
Hanna Borlinghaus & \emph{University of Konstanz, Germany}\\
Emek Demir    & \emph{OHSU, USA}\\
Andreas Dräger   & \emph{University of Tübingen, Germany}\\
Robin Haw   & \emph{Ontario Institute for Cancer Research, Canada}\\
Alexander Mazein   & \emph{University of Luxembourg, Luxembourg}\\
Huaiyu Mi	& \emph{University of Southern California, USA}\\
Stuart Moodie   & \emph{Eight Pillars Ltd, UK}\\
Falk Schreiber	 & \emph{University of Konstanz, Germany}\\
Anatoly Sorokin   & \emph{Institute of Cell Biophysics RAS, RU}\\ 
Vasundra Touré   & \emph{NTNU, Norway}\\
Alice Vill\'{e}ger   & \emph{Freelance IT Consultant, UK}\\
Tobias Czauderna   & \emph{University of Applied Sciences Mittweida, Germany}\\
Ugur Dogrusoz   & \emph{Bilkent University, Turkey}\\
Augustin Luna   & \emph{NIH, USA}\\
\end{tabular}

\vfill

\normalsize
\begin{minipage}{5in}
  \emph{To discuss any aspect of SBGN, please send your messages
    to the mailing list \mailto{sbgn-discuss@googlegroups.com}.  To get
    subscribed to the mailing list or to contact us directly,
    please write to \mailto{sbgn-editors@googlegroups.com}. Bug reports and specific comments about the specification should be entered in the issue tracker \url{https://github.com/sbgn/process-descriptions/issues}.}
\end{minipage}

\vfill

\centerline{\includegraphics[width=1.25in]{images/build/sbgn_logo.pdf}}

\end{center}

\end{titlepage}

% The title page is considered unnumbered and the next page after this
% starts with the page number 1 (actually, i), but the duplication of page
% number 1 confuses hyperref.
% The following change makes the title page have page number 1 and the next
% page after that it becomes page ii.  This is unorthodox, but seems not
% completely unreasonable, and it avoids the confusing hyperref.

\setcounter{page}{2}
